% chapters/01_introduction_and_terminologies.tex
% 第1章：引言与术语
\chapter{引言：重新思考“训练”与“测试”的边界}
\label{ch:intro}

\epigraph{``The dichotomy between training and testing is an artificial artifact of how we build algorithms, not a reflection of how intelligent systems actually operate in the world.''}{--- 引言}

\section{本章想回答什么问题}

在传统的深度学习教育中，我们被灌输了一个极其极其根深蒂固的概念：\textbf{训练与测试是相互独立的两个阶段}。
\begin{itemize}
    \item \textbf{训练阶段}：需要大量的数据、GPU 集群、耗时数周的反向传播计算。这时模型的参数 $\theta$ 在不断变化。
    \item \textbf{测试/推理阶段}：模型被冻结了，参数 $\theta$ 锁定不变。模型被发往生产环境执行前向预测。不管在这个阶段模型犯了多少荒谬的错误，或者见到了多少明显的新规律，它都不会有丝毫长进。
\end{itemize}

然而，真实世界的数据分布永远在流动（Distribution Shift）。这种被冻结的“死”模型常常会在遇到一点点没有见过的输入时全盘崩溃。
本章，乃至本书想要回答的一个核心问题是：
\begin{quote}
\textit{为什么我们在测试（推断）时不能也让模型继续学习和适应？在极端的算力与架构重塑下，我们能否打破这一道被视为经典的铁丝网，让“测试时”本身就是“训练时”？}
\end{quote}

我们将从基础的术语定义开始，引入本书反复提及的基石问题与数学环境，并帮助读者在脑海中建立对动态自适应系统的生物学直觉。

\section{最小例子：面对不断变轨的小球}
\label{sec:intro_minimal}

\begin{toybox}[title=玩具问题 A：预测弹跳路径分布的偏移]
假设我们训练了一个简单的模型来预测桌球的弹跳轨迹（回归问题）。在训练集中，桌台的摩擦系数是确定的 $f=0.1$。模型在预测试阶段完美学会了给出下三个时间点的位置 $(x,y)$。

\textbf{部署（测试）时：}小桌台的布面被替换了，摩擦系数隐性变成了 $f=0.5$（发生了分布偏移 Distribution Shift）。
\begin{itemize}
    \item \textbf{传统僵化参数模型}：看到球弹出后，仍旧按照 $f=0.1$ 的物理公式去推断，每次都会预测小球滑得更远。尽管模型“看到”了球实际上停在更近的地方，它不具备参数更改能力，下一次依然预测错误。
    \item \textbf{测试时适应系统（Test-Time Adaptive System）}：在第一根球的轨迹发生后，它将\textbf{本次不完美的输出预测和观测的偏差}作为短效损失，在微秒内部就对模型做了一次瞬时的浅层（或隐式）再调整。到第三根球时，它已经隐约知道需要“把预测缩短”，从而适应了新桌布。
\end{itemize}
这就是在测试时实施学习的最原始直觉展现。
\end{toybox}

\section{直觉解释：静态字典 vs. 活着的大脑}
\label{sec:intro_intuition}

要理解后续将要介绍的所有前沿架构，首先我们需要剥离对“参数就是绝对真理”的执念。

在常规范式下，我们训练大语言模型就好比是在印刷一本包罗万象的厚重大字典（预训练）。字典印完之后（进入测试），你来查询任何东西，它都是以固定页面结构直接反射回去的。

而测试时学习的思想是：预训练不应当仅仅学习“字典固有的页面内容（参数）”，更多的是要学习去搭建一个“能在被提问时立刻用身边的草稿纸结合你的上文去打个草稿解题的小助理机制”。你需要的不是去读一本死书，而是在提供一种可供在线校准的思考架构引擎。

\begin{neurosciencebox}[title=神经科学直觉：选择性注意与快速适应]
在神经科学中，这种现象非常常见。即使是脑部已经完成了“基本常识结构发育”（长期预训练），人类在进入一个全然陌生的暗室时：
\begin{itemize}
    \item \textbf{计算角色上的相似：}由于光线的改变，人的视觉皮层并不需要回到“娘胎里”（重新进入漫长的预训练中心）重组架构，而是在最初的几秒适应内利用局部的误差预期，迅速改变当前处理通道的敏感电位（在线快速适应）。
    \item \textbf{工作记忆分配：}遇到长上下文的理解任务（例如听长篇演讲），大脑的前额叶皮层拥有选择性注意（Selective Attention）能够临时把某一句重要的话“贴在脑前”（推断阶段的局部存储与计算），而非把它写死到长达十年寿命的长期记忆神经元里去。
\end{itemize}
测试时间的学习，本质就是在计算机工程中去唤醒这种本该存在的“现场试错适应”（Sensorimotor calibration）以及短时间尺度的工作缓存运算能力。
\end{neurosciencebox}

\begin{analogyboundarybox}[title=这并不是说“模型就是大脑”]
虽然从直觉上我们可以把“在测试环节自我纠偏”类比为大脑的临场适应反应，但你绝不能因此建立以下推论：
\begin{itemize}
    \item \textbf{实现机制上的云泥之别：}人工模型依赖的是计算出精确的标量损失，进而利用明确的数学雅可比求导链甚至二阶反置矩阵实施“向回退”的精确修正同步。真实大脑在类似短时调整中不依靠标准的反向传播，而是大量利用局部电位扰动和多模式神经递质（像是多巴胺漫灌带来的时标突触可变）发生。
    \item 我们在人工智能里探讨的设计是对人类这种“具有多个反馈速度环”功能的\textbf{数学粗糙模拟仿化}，而不是对大脑生理器官的数字重现！它只是在“功能组织目的性”上有高度一致，但在数学解剖下是两种截然不同的演化路线。
\end{itemize}
\end{analogyboundarybox}

\section{数学形式化与核心术语}
\label{sec:intro_math}

为了整本书的讨论同出一辙，这里严格定义我们的基础场景与数学标记。

\subsection{时间轴的区分}
\begin{itemize}
    \item $t_\text{train}$：指代主模型被喂养海量人类通用数据集以定型宏观静态参数的过程。（俗称：外循环，预训练）。
    \item $t_\text{test}$ 或 $t$：模型被部署，面对特定任务时的单个序列推移阶段（俗称：内循环，推断时间步骤）。
\end{itemize}

\subsection{参数空间的切割}
我们将参数概念做了彻底解绑：
\begin{itemize}
    \item \textbf{慢权重（Slow Weights / Global Parameters）} $\theta$：在 $t_\text{train}$ 通过 SGD/Adam 从巨大训练集中极其缓慢地固定下来。在推演 $t_\text{test}$ 时理论上被冻结了不会发生显存重写式变化，它们承载\textbf{全局常识与提取规则}。
    \item \textbf{快权重 / 隐状态参数（Fast Weights / Inner State）} $\mS_t$：仅在这个特定序列推演的 $t_\text{test}$ 开始时由零初构，在推算过程的毫秒间隙受输入数据影响实时疯狂改写。在整个推理序列结束时被直接从内存抛弃删除！它们承载\textbf{本地适应与特定任务偏好}。
\end{itemize}

\subsection{普通推断 vs. 自适应推断的基础方程异同}
普通序列推断的函数定义为在给定慢权重 $\theta$ 后：
\begin{equation}
  y_t = f(x_t; \theta, \mS_{t-1}) \quad \text{其中} \quad \mS_{t} = g(x_t, \mS_{t-1}; \theta)
\end{equation}
在这里隐藏状态 $\mS$ 的更新 $g$ 是一个在设定死的一组逻辑下的普通变换组合（如加减拼接）。

而测试时学习体系下：
\begin{equation}
   \mS_t \leftarrow \text{Optimizer}_\eta(\mS_{t-1}, \mathcal{L}_{\text{inner}}(x_t; \mS_{t-1}, \theta))
\end{equation}
这里的 $\mS_{t-1}$ 被视为一组被接受优化迭代的小型网络内部参数组。$\mathcal{L}_{\text{inner}}$ 被定义为某个无监督或是由于上下文自然生成出的误差拟合源项！注意：它是求极值反向退算的优化过程，而不是固定前向！

\section{两条 TTT 线索的本质差异}

\subsection*{第一条线索：TTT 2020 风格（外科手术式微调）}

\citet{sun2020ttt} 提出的"Test-Time Training" 指的是：
当测试样本 $x_\text{test}$ 到来时，利用一个\textbf{自监督辅助任务}
对网络参数 $\theta$ 执行真实梯度下降更新，然后再做主任务预测。

\begin{center}
\texttt{输入测试样本} $\to$ \texttt{自监督微调（修改 $\theta$）} $\to$ \texttt{主任务预测}
\end{center}

关键点：这需要调用标准的反向传播（\texttt{.backward()}）。
每处理一个测试样本，都要运行一次或多次梯度步骤。
这种方法\textbf{真实地改变了模型的持久化参数}。这就是第 2 层更新。

\subsection*{第二条线索：TTT-Layer 2024 风格（内嵌计算式更新）}

\citet{sun2024ttt} 以及相关工作（见第 \ref{ch:ttt_layers} 章）
把"TTT" 重新定义为：
将序列模型的\textbf{隐藏状态} $s_t$ 本身看作一个微型学习器的参数，
在前向传播过程中，每个新 token 到来时就执行一次基于状态的梯度步骤。

\begin{center}
\texttt{前一隐藏状态 $s_{t-1}$} $\to$ \texttt{输入 $x_t$} $\to$
\texttt{内部梯度步骤（更新 $s_t$）} $\to$ \texttt{输出 $y_t$}
\end{center}

关键点：\emph{整个过程都在前向传播内完成}。
没有外部的 \texttt{.backward()} 调用。
模型的\textbf{持久化参数 $\theta$ 不变}；变化的是每次推断的隐状态。这就是第 1 层更新。

\section{代表文献指路}
\label{sec:intro_papers}

由于本章为总计概念梳理引言。我们将本教材引用的整个领域发展线，切分为三部曲，后附关键奠基文章：
\begin{enumerate}
    \item \textbf{概念开局与人工修剪期}（涉及第2,3,4章）：在发现普通模型极其脆弱后，研究者直接在推断流程前用代码加入一个小的微调。\citet{sun2020ttt}《Test-time training with self-supervision for generalization under distribution shifts》是这一领域的正式定命之作。
    \item \textbf{架构内化隐式化期}（涉及第5,6章）：探索为什么无需人工手改代码而仅凭多层大结构的自然前向传递本身，竟然神奇地实现了类似在线训练的效果（即ICL本质证明），如 \citet{oswald2022transformers}；以及使用快权重代替显式反算的起步。
    \item \textbf{全面内部结构重铸期}（涉及第 7 至 11 章）：将大模型底层进行大修，把 RNN、注意力全部整合为“自包含隐状态回归层”，如 \citet{sun2024ttt} 的 TTT-layers，直到向终极嵌套多阶系统演化迈进。
\end{enumerate}

\section{和后序章节的关系}
\label{sec:intro_bridges}

本章负责确立马步桩台。
理解了“$\theta$ 产生寻找适应梯度的能力”和“$\mS_t$ 被适应梯度去改写自身”的双元区别之后，紧接的第 \ref{ch:why_test_time} 章就会直接从经典的计算机视觉例子去直观解释当分布（比如下雨天的模糊图像导致与训练严重不符）开始脱轨时，我们在测试点现场强行加几步优化梯度的具体实施方式以及其力挽狂澜的作用。

\begin{labbox}[title=Lab 1: 验证静态参数大模型在简单位移下有多“脆弱”]
\textbf{目标：}手动编写脚本并观察纯固定参数在遭遇 OOD (Out-of-Distribution) 未见系统时的预测灾难。并实施一次最简单的“测试补偿”。

\begin{enumerate}
    \item 构建一个两层的全连接神经网络用于给二维坐标打点做分类。分布是 $y = 2x$ 这条斜线上下的分离。
    \item \textbf{训练期：}数据全为 $[-1, 1]$ 内取样，使用 PyTorch 训练直到损失$<0.01$；\textbf{切断}优化器更新机制（\texttt{eval()}态）。
    \item \textbf{测试期：}开始送入 $+50$ 整体位移后的 $[49, 51]$ 特征域空间的新样本点。
    \item \textbf{观察：}虽然线性关联未变，但是由于未见过高分位数范围值，静态的分类网络前向非线性激活发生过激失真，原本完全正确的判断预测输出全部降到了近乎 $50\%$ 的瞎猜状态。
    \item \textbf{补偿（实验核心）：}如果在部署阶段，我们将这批新数据自己和自己构成的均值或进行一个简单的自适应自编码重构误差（人工给一个无监督 Loss），然后仅对预测层前进行 \textbf{1次 SGD 现场反向推正}。再去分类打标。你会发现在毫秒之间，精度又拉回了 $98\%$。
\end{enumerate}
这就是本书致力于追求去系统化的、发生于推测期的核心魅力。
\end{labbox}

\begin{misconceptionbox}
\textbf{常见误解 1：测试时的训练是在给大字典加加页数扩容。}

\textbf{纠正：}不！这是典型的死记硬背理解。模型在推断时的学习操作更多地是为了\textbf{提取关联与抑制主预训练库中的无关特征}，是将特定数据的模式（比如当前文章的写作风格）暂时置顶提升的在线优化工作，它处理过后立刻就会随着当前提问消失。这绝对不是说把这次用户的提问永久固化进去了大模型的基座文件池。
\end{misconceptionbox}

\begin{misconceptionbox}
\textbf{常见误解 2：我们现在常用的 ChatGPT 就是拥有这种不断根据我的对话进行内部极速梯度改写本领才看起来这么聪明的。}

\textbf{纠正：}在撰写这本教材的绝大多数前线工业模型（GPT-4时代等）并\textbf{没有}采用内化的测试梯度层（TTT系列还在学术至工业的飞速跨越期中）。大多数线上产品能针对你的语境显得“能自我调节”，全靠极其暴力的 Transformer 并发超长缓存检索历史序列硬生生建立多维组合（即第 6 章讲的隐式 ICL），它们并不是去更改什么底层态参数本身进行的，这也正是极其昂贵并在将来必定需要通过“参数层即是学习者”演变所必然取代的源头。
\end{misconceptionbox}

\section{练习题}
\label{sec:intro_exercises}

\begin{exercise}
\textbf{概念对位理解：}请用三句话分别解释以下概念并且指出它们产生及作用的阶段：(1) 全局慢权重 $\theta$；(2) 分布偏移 (Distribution Shift)；(3) 工作记忆与隐藏态学习器。并明确标出哪一个概念更贴近于在传统预训练（Pre-train）中所侧重的。
\end{exercise}

\begin{exercise}
\textbf{数学定义的扩展：}如果在上面数学形式化章节中，针对公式 $g$ 更新隐态所使用的损失函数 $\mathcal{L}_{\text{inner}}$ 被彻底置零了，上述方程会不可避免地退化为什么？而此时如果还需要根据输入数据维持不同的回答可能性输出，只能借由让网络外部变态级加大哪一项去进行对过去的死记硬背？这体现了什么代价的权衡（Trade-off）？
\end{exercise}

\begin{exercise}
\textbf{关于真实应用成本：}如果我们如上述玩具问题中的小球或者位移打标情况，每一次都必须要人工去计算一种叫做自恢复误差来进行推断时候的反向传播以获得能力。它在实际手机端部署上的最大两项无法接受的阻碍会是什么（结合框架环境和计算图建立考量）？
\end{exercise}

\section{本章小结}
\label{sec:intro_summary}

\begin{itemize}
    \item \textbf{划开坚冰壁垒：}打破“训练完后就只管输出永远不再接受改写”的静态推算理念！提出了需要且必需针对输入流发生不可预测移位时执行实时、小规模优化的全新宏大议程。
    \item \textbf{确立概念范式：}给出了主从区分：在主端（具有几近不变、承载泛化规则提炼特征的网络骨架的全局慢权重）与局侧前线（面对短序列新数据流立刻展开梯度回归学习改写的极速抛弃式隐态）。
    \item \textbf{提供了研究锚点：}这本教材不仅仅是写给研究人员的工具合集，更是为了改变你认知深度网络“智能构成成分结构”底层思想观的思维体操。
\end{itemize}

\paragraph{读完本章，你应该能够：}
\begin{itemize}
    \item 准确抛弃“训练和测试截然不可并协互通”的陈规旧俗。
    \item 采用“系统预先学会一种打草稿寻找短时修正目标的能力”的视角看待为何大模型不会完全在测试态卡死。
    \item 掌握全局结构慢权重 $\theta$ 与 伴数据流局部发生作用变造的隐藏快状态 $\mS_t$ 此一通贯全书的双轨分类学术语标定。
\end{itemize}
